\chapter{L0 — Recursion-theoretic prerequisites and Proposition B}

% P-R reference: §0.2 "Prerequisites from logic and analysis",
% literature/papers/PourEl-Richards-0.2.prerequisites.md.

\section{Recursion-theoretic bridge}

P-R uses only two nontrivial facts from recursion theory: the existence of a
recursively enumerable nonrecursive set (Proposition A, this section) and the
existence of a recursively inseparable pair (Proposition B, next section). Both
are stated against Mathlib's \texttt{Computability.Primrec} / \texttt{Partrec}
machinery rather than against the partial-recursive grammar P-R describes
informally.

\begin{definition}
  \label{def:isRecursive}
  \lean{ComputableAnalysis.L0.IsRecursive, ComputableAnalysis.L0.IsRecursivelyEnumerable, ComputableAnalysis.L0.IsRecursiveSet}
  \leanok
  For \texttt{Primcodable} types $\alpha, \beta$:
  \begin{itemize}
    \item \textbf{IsRecursive} $f$ abbreviates \texttt{Computable}~$f$ (a total recursive function $\alpha \to \beta$).
    \item \textbf{IsRecursivelyEnumerable} $A$ abbreviates \texttt{REPred} $(\lambda x.\,x \in A)$ (semi-decidable membership).
    \item \textbf{IsRecursiveSet} $A$ abbreviates \texttt{ComputablePred} $(\lambda x.\,x \in A)$ (decidable membership).
  \end{itemize}
  These are pure aliases of Mathlib symbols introduced to keep P-R Ch.\,0.2 vocabulary visible inside the project.
\end{definition}

\begin{definition}
  \label{def:cantorPair}
  \lean{ComputableAnalysis.L0.cantorPair}
  \leanok
  The Cantor pairing function \texttt{cantorPair}\,$: \mathbb{N} \to \mathbb{N} \to \mathbb{N}$, an alias of Mathlib's \texttt{Nat.pair}. P-R Ch.\,0.2 implicitly assumes a primitive-recursive bijection $\mathbb{N}^2 \to \mathbb{N}$; \texttt{Nat.pair} is the canonical choice.
\end{definition}

\begin{lemma}
  \label{lem:cantorPair_computable}
  \lean{ComputableAnalysis.L0.cantorPair_computable}
  \leanok
  \uses{def:cantorPair}
  \texttt{cantorPair} is \texttt{Computable}$_2$.
\end{lemma}
\begin{proof}\leanok
  \texttt{Nat.pair} is built from projections and arithmetic that are all
  primitive recursive, hence computable; the binary form follows by lifting
  \texttt{Primrec}~$\to$~\texttt{Computable}.
\end{proof}

\begin{definition}
  \label{def:haltingSetAt}
  \lean{ComputableAnalysis.L0.haltingSetAt}
  \leanok
  For a fixed input $n \in \mathbb{N}$, \texttt{haltingSetAt}~$n$ is the set of \texttt{Nat.Partrec.Code} indices whose interpretation halts on input $n$. This is Mathlib's $K_n$-style halting set at $n$.
\end{definition}

\begin{lemma}
  \label{lem:haltingSetAt_re}
  \lean{ComputableAnalysis.L0.haltingSetAt_re}
  \leanok
  \uses{def:isRecursive, def:haltingSetAt}
  \texttt{haltingSetAt}~$n$ is recursively enumerable.
\end{lemma}
\begin{proof}\leanok
  Membership of a code $c$ is semi-decidable: evaluate $c$ on input $n$ and
  accept exactly when the computation halts. This is the acceptance domain of a
  partial-recursive predicate, so the set is r.e.
\end{proof}

\begin{lemma}
  \label{lem:haltingSetAt_not_recursive}
  \lean{ComputableAnalysis.L0.haltingSetAt_not_recursive}
  \leanok
  \uses{def:isRecursive, def:haltingSetAt}
  \texttt{haltingSetAt}~$n$ is not recursive.
\end{lemma}
\begin{proof}\leanok
  A decision procedure for membership in \texttt{haltingSetAt}~$n$ would decide
  whether an arbitrary code halts at $n$, i.e.\ solve the halting problem, which
  is impossible.
\end{proof}

\begin{theorem}[Proposition A]
  \label{thm:propA}
  \lean{ComputableAnalysis.L0.prop_A_exists_re_not_recursive}
  \leanok
  There exists a recursively enumerable set $A \subseteq \mathbb{N}$ that is not recursive. P-R Ch.\,0.2:12.
\end{theorem}
\begin{proof}\leanok
  \uses{lem:haltingSetAt_re, lem:haltingSetAt_not_recursive}
  Take $A = \texttt{haltingSetAt}~n$: it is recursively enumerable by
  \Cref{lem:haltingSetAt_re} and not recursive by
  \Cref{lem:haltingSetAt_not_recursive}.
\end{proof}

\section{Proposition B — a recursively inseparable r.e. pair}

\begin{definition}
  \label{def:insepAB}
  \lean{ComputableAnalysis.L0.insepA, ComputableAnalysis.L0.insepB}
  \leanok
  \texttt{insepA} and \texttt{insepB} are subsets of \texttt{Nat.Partrec.Code}: \texttt{insepA} is the set of codes that, applied to their own index, halt with output $0$; \texttt{insepB} the same but with output $1$. P-R Ch.\,0.2:35.
\end{definition}

\begin{lemma}
  \label{lem:insep_re}
  \lean{ComputableAnalysis.L0.insepA_re, ComputableAnalysis.L0.insepB_re}
  \leanok
  \uses{def:insepAB, def:isRecursive}
  Both \texttt{insepA} and \texttt{insepB} are recursively enumerable.
\end{lemma}
\begin{proof}\leanok
  Each set is the acceptance domain of a partial-recursive predicate: run a code
  on its own index and accept on output $0$ (resp.\ $1$). Semi-decidability of
  that predicate gives recursive enumerability.
\end{proof}

\begin{lemma}
  \label{lem:insep_disjoint}
  \lean{ComputableAnalysis.L0.insepA_insepB_disjoint}
  \leanok
  \uses{def:insepAB}
  \texttt{insepA} $\cap$ \texttt{insepB} $= \emptyset$.
\end{lemma}
\begin{proof}\leanok
  A single computation cannot return both $0$ and $1$, so no code lies in both
  sets.
\end{proof}

\begin{theorem}
  \label{thm:insep_no_separator}
  \lean{ComputableAnalysis.L0.insepA_insepB_no_separator}
  \leanok
  \uses{def:insepAB, def:isRecursive}
  No recursive set $C$ satisfies \texttt{insepA} $\subseteq C$ and $C \cap$ \texttt{insepB} $= \emptyset$.
\end{theorem}
\begin{proof}\leanok
  Suppose such a recursive $C$ existed. Apply the recursion theorem
  (\texttt{Nat.Partrec.Code.fixed\_point}$_2$) to the characteristic function of
  $C$ to obtain a self-referential code that halts with output $1$ when its own
  index lies in $C$ and output $0$ otherwise. That code is then forced into
  \texttt{insepB} when in $C$ and into \texttt{insepA} when outside $C$,
  contradicting $\texttt{insepA} \subseteq C$ and $C \cap \texttt{insepB} = \emptyset$.
\end{proof}

\begin{theorem}[Proposition B]
  \label{thm:propB}
  \lean{ComputableAnalysis.L0.prop_B}
  \leanok
  There exist disjoint recursively enumerable sets $A, B \subseteq \mathbb{N}$ such that no recursive set separates them. P-R Ch.\,0.2:35.
\end{theorem}
\begin{proof}\leanok
  \uses{lem:insep_re, lem:insep_disjoint, thm:insep_no_separator}
  Take $A = \texttt{insepA}$ and $B = \texttt{insepB}$: both are r.e.\
  (\Cref{lem:insep_re}), disjoint (\Cref{lem:insep_disjoint}), and admit no
  recursive separator (\Cref{thm:insep_no_separator}).
\end{proof}

\section{Analysis prerequisites pointer}

\begin{remark}
  \label{rem:analysisBridge}
  The file \texttt{ComputableAnalysis/L0/AnalysisBridge.lean} records the
  P-R\,$\leftrightarrow$\,Mathlib correspondence for the analysis hierarchy
  (\texttt{NormedSpace}, \texttt{InnerProductSpace}, \texttt{lp},
  \texttt{MeasureTheory.Lp}, \texttt{ContinuousMap}, \texttt{ContDiff}). It
  introduces no symbols: downstream layers (L3, L4) import Mathlib's names
  directly. See P-R Ch.\,0.2:58--81 for the analysis prerequisites P-R lists.
\end{remark}
